\documentclass[12pt,a4paper]{report}
\usepackage[utf8]{inputenc}
\usepackage[T1]{fontenc}
\usepackage[english]{babel}
\usepackage{amsmath, amssymb, amsthm, amsfonts}
\usepackage{geometry}
\usepackage{xcolor}
\usepackage{listings}
\usepackage{tikz}
\geometry{hmargin=2.5cm,vmargin=2.5cm}

% Theorem-like environments
\newtheorem{theorem}{Theorem}[section]
\newtheorem{proposition}[theorem]{Proposition}
\newtheorem{lemma}[theorem]{Lemma}
\newtheorem{corollary}[theorem]{Corollary}
\theoremstyle{definition}
\newtheorem{definition}{Definition}[section]
\newtheorem{example}{Example}[section]
\theoremstyle{remark}
\newtheorem{remark}{Remark}[section]
\newtheorem{property}{Property}[section]

% Python listings
\lstset{
    language=Python,
    basicstyle=\ttfamily\small,
    keywordstyle=\color{blue}\bfseries,
    stringstyle=\color{red},
    commentstyle=\color{green!60!black}\itshape,
    numbers=left,
    numberstyle=\tiny\color{gray},
    stepnumber=1,
    frame=single,
    breaklines=true
}

\title{Fundamentals of Natural Language Processing\\[0.3cm] \large Chapter 2 --- Text Representation and Evaluation}
\author{Aymen Ben Brik}
\date{}

\begin{document}

\maketitle

\setcounter{chapter}{1}
\chapter{Text Representation and Evaluation}

A computer cannot operate on raw text. Before any model can ``read'' a sentence, the sentence must be \emph{tokenized} (cut into discrete units), \emph{numericalized} (converted into integer indices), and \emph{vectorized} (embedded in a vector space). This chapter develops the chain rigorously: tokenization granularities, sparse and dense vector-space models, semantic similarity, and the evaluation metrics by which we compare NLP systems.

\section{Tokenization fundamentals}

\subsection{Motivation}
Tokenization is the unsung hero of every NLP pipeline. A poor tokenizer silently undermines the most sophisticated downstream model: out-of-vocabulary words become unknown tokens, morphological variants are treated as unrelated, multi-word expressions are fragmented. Conversely, a tokenizer aligned with the data distribution can shrink vocabulary size, accelerate training, and improve generalisation. Modern Large Language Models all rely on a single technical choice --- \emph{sub-word tokenization} --- whose strengths and limitations every practitioner should understand.

\subsection{Approach by problem: how many tokens are in this sentence?}
\textbf{Problem.} Consider the sentence
\begin{quote}
\emph{``I've been to Sfax-University; it's amazing!''}
\end{quote}
Count how many ``tokens'' it contains under each of the following naive policies:
\begin{enumerate}
    \item split on whitespace;
    \item split on whitespace and punctuation;
    \item split into individual characters.
\end{enumerate}

\textbf{Solution sketch.}
\begin{enumerate}
    \item Whitespace only: $\{$\texttt{I've}, \texttt{been}, \texttt{to}, \texttt{Sfax-University;}, \texttt{it's}, \texttt{amazing!}$\}$ --- $6$ tokens. Punctuation glued to words; ``I've'' is a single token.
    \item Whitespace and punctuation: $\{$\texttt{I}, \texttt{'ve}, \texttt{been}, \texttt{to}, \texttt{Sfax}, \texttt{-}, \texttt{University}, \texttt{;}, \texttt{it}, \texttt{'s}, \texttt{amazing}, \texttt{!}$\}$ --- $12$ tokens. The contraction \emph{I've} is split, but so is the legitimate compound \emph{Sfax-University}.
    \item Characters: $35$ tokens. Punctuation, capitalisation and spaces are all preserved, but every word is fragmented; the model would have to relearn ``what is a word'' from data.
\end{enumerate}
Each policy entails a trade-off between vocabulary size and information preserved. There is no universally correct answer; the choice is engineering, guided by the downstream task and the language being modelled.

\subsection{Definition and the tokenization pipeline}

\begin{definition}[Tokenization]
\textbf{Tokenization} is the process of converting a raw text string into an ordered sequence of discrete units called \textbf{tokens}, drawn from a finite inventory called the \textbf{vocabulary} $V$.
\end{definition}

\begin{definition}[Numericalization]
\textbf{Numericalization} is the bijective mapping $V \leftrightarrow \{0, 1, \dots, |V|-1\}$ that assigns each token an integer identifier.
\end{definition}

The full pipeline is: \emph{raw text} $\to$ tokens $\to$ token IDs $\to$ tensor inputs to a model. Every NLP system is built on top of these three trivially simple but practically critical steps.

\subsection{Three levels of granularity}

\begin{center}
\begin{tabular}{|l|l|l|l|}
\hline
\textbf{Level} & \textbf{Vocabulary size} & \textbf{Sequence length} & \textbf{OOV behaviour} \\
\hline
Word & Very large ($10^5$--$10^6$) & Short & Unknown words become \texttt{[UNK]} \\
Character & Tiny ($\sim 100$) & Very long & No OOV (any character is in $V$) \\
Sub-word & Moderate ($10^4$--$10^5$) & Moderate & Robust: any word splits into known pieces \\
\hline
\end{tabular}
\end{center}

\paragraph{Word tokenization.}
The intuitive default. Tokens are entire words separated by whitespace and punctuation rules. Strengths: each token carries lexical meaning, and downstream models converge faster on a small annotated dataset. Weaknesses: vocabulary explodes for morphologically rich languages (\emph{Arabic, Turkish, Finnish}); typos and rare words become \texttt{[UNK]}; cannot generalise across morphological variants (\emph{run / runs / running} are three unrelated tokens).

\paragraph{Character tokenization.}
Each character is a token. Strengths: vocabulary is bounded and stable across languages; no out-of-vocabulary problem (any unseen word becomes a sequence of known characters). Weaknesses: sequences become $5$--$10\times$ longer, increasing both memory and compute; the model must learn from scratch which character sequences form meaningful units.

\paragraph{Sub-word tokenization.}
The dominant choice in 2020s NLP. Vocabulary contains \emph{frequent words as wholes} and \emph{rare words split into meaningful pieces}. The split is learnt from data using algorithms such as Byte-Pair Encoding (BPE), WordPiece (BERT) and SentencePiece (XLM-R, T5). Sub-word tokenization combines the strengths of word-level (semantic units for common words) and character-level (no OOV) approaches.

\subsection{Pre-processing operations}

Tokenization is often combined with ad-hoc transformations that modify or normalise the input string before splitting.

\begin{itemize}
    \item \textbf{Lowercasing}: ``Apple'' and ``apple'' become equivalent. Useful for spelling-insensitive tasks but destroys signal in NER (a capitalised \emph{Apple} is more likely the company).
    \item \textbf{Punctuation removal}: simplifies vocabulary but loses sentence boundaries.
    \item \textbf{Stop-word removal}: discard high-frequency words (\emph{the, of, in}) that carry little task-specific information. Useful for keyword search; harmful for syntax-sensitive tasks.
    \item \textbf{Stemming}: chop morphological suffixes by hand-coded rules (Porter, Snowball). \emph{running} $\to$ \emph{run}. Fast but linguistically crude.
    \item \textbf{Lemmatization}: map a word to its dictionary form using a morphological analyser. \emph{better} $\to$ \emph{good}. Slower but linguistically principled.
\end{itemize}

\begin{remark}
Modern large language models perform little or no manual pre-processing: they are trained on raw text with sub-word tokenization, and the model itself learns whatever normalisations it needs. Heavy pre-processing is typical of statistical-era pipelines and remains useful for small-data, classical-ML systems.
\end{remark}

\subsection{Byte-Pair Encoding (BPE) in detail}

BPE is the simplest sub-word algorithm and the workhorse of GPT-2, GPT-3, GPT-4, RoBERTa, and many others. It iteratively merges the most frequent pair of adjacent symbols.

\paragraph{Algorithm.}
\begin{enumerate}
    \item Start with a vocabulary that contains every individual character of the corpus.
    \item Treat each word as a sequence of characters with an end-of-word marker.
    \item Count the frequency of every adjacent symbol pair across the corpus.
    \item Merge the most frequent pair into a new symbol; add it to the vocabulary.
    \item Repeat steps 3--4 until a target vocabulary size is reached.
\end{enumerate}

\begin{example}[A miniature BPE run]
Suppose the corpus consists of the words \texttt{low}, \texttt{lower}, \texttt{newest}, \texttt{widest}, with frequencies $5, 2, 6, 3$. The initial vocabulary is $\{$\texttt{l, o, w, e, r, n, s, t, i, d}, \texttt{</w>}$\}$. Counting pairs:
\[ \texttt{e\,s} = 6 + 3 = 9,\quad \texttt{s\,t} = 9,\quad \texttt{l\,o} = 7, \quad \texttt{o\,w} = 7, \dots \]
The pair \texttt{e\,s} (9 occurrences) is merged first into a new symbol \texttt{es}. Recount, merge again --- now \texttt{es\,t} (9 occurrences) into \texttt{est}. Continue: \texttt{lo}, \texttt{low}, \texttt{ne}, \texttt{new}, \dots\ After enough merges, frequent words like \texttt{low</w>} become single tokens, while rare words remain decomposed. A new word like \texttt{lowest} is then tokenized as \texttt{low + est}, even though it never appeared in training.
\end{example}

\paragraph{Why BPE works.} The algorithm provides a natural rate--distortion trade-off. As vocabulary grows, the average sequence length shrinks, but the ``unit cost'' per token increases (each token carries more information). Empirically, vocabularies of $30$k--$100$k sub-words give an excellent trade-off across many languages.

\subsection{Multilingual and dialectal challenges}

Tokenization choices that are adequate for English may fail for other languages.

\begin{itemize}
    \item \textbf{Whitespace-free scripts.} Chinese, Japanese, Thai write without spaces. Word tokenization requires segmentation algorithms (e.g.\ Jieba, MeCab).
    \item \textbf{Morphologically rich languages.} Arabic and Turkish concatenate prefixes and suffixes onto a stem, generating very large surface-form inventories. Sub-word tokenization is essential.
    \item \textbf{Code-switching and dialects.} Tunisian Arabic frequently mixes Modern Standard Arabic, French, and Latin-script transliterations (``arabizi''). A general-purpose multilingual tokenizer (e.g.\ XLM-R's SentencePiece) is recommended over single-language tokenizers.
    \item \textbf{Right-to-left scripts.} Display order and logical order can differ; numerical IDs operate on logical order, but evaluation tools must be aware of the visual rendering.
\end{itemize}

\subsection{Worked examples}

\textbf{Example 1 (Granularity choice --- Easy).}
You are training a classifier on $1000$ short SMS messages, mostly in English with frequent typos. Which tokenization granularity is most appropriate, and why?

\textit{Solution.} Sub-word tokenization. Word-level would suffer from OOV (typos produce never-seen surface forms); character-level would lengthen sequences and dilute the signal in a small dataset. Sub-word splits typos into recognised pieces (\texttt{hapy} $\to$ \texttt{ha + py}) without exploding the vocabulary.

\medskip

\textbf{Example 2 (BPE iteration --- Medium).}
Apply two iterations of BPE to the toy corpus $\{\texttt{aaab}, \texttt{aabb}, \texttt{abab}\}$ (each of frequency $1$). Show the resulting vocabulary after each merge.

\textit{Solution.} Initial vocabulary: $\{\texttt{a}, \texttt{b}\}$.

Iteration 1: count pairs --- \texttt{a\,a} appears in \texttt{aaab} (twice) and \texttt{aabb} (once) and \texttt{abab} (zero) $= 3$; \texttt{a\,b} appears in \texttt{aaab} (once) and \texttt{aabb} (once) and \texttt{abab} (twice) $= 4$; \texttt{b\,a} appears once (in \texttt{abab}); \texttt{b\,b} once.
Most frequent: \texttt{a\,b} (4). Merge into \texttt{ab}. Vocabulary: $\{\texttt{a}, \texttt{b}, \texttt{ab}\}$. Re-tokenization: \texttt{aaab} $\to$ \texttt{a\,a\,ab}; \texttt{aabb} $\to$ \texttt{a\,ab\,b}; \texttt{abab} $\to$ \texttt{ab\,ab}.

Iteration 2: count new pairs --- \texttt{a\,a} (in \texttt{aaab}, once) $= 1$; \texttt{a\,ab} (in \texttt{aaab} and \texttt{aabb}) $= 2$; \texttt{ab\,b} (in \texttt{aabb}) $= 1$; \texttt{ab\,ab} (in \texttt{abab}) $= 1$.
Most frequent: \texttt{a\,ab} (2). Merge into \texttt{aab}. Vocabulary: $\{\texttt{a}, \texttt{b}, \texttt{ab}, \texttt{aab}\}$.

\medskip

\textbf{Example 3 (Tokenizer mismatch --- Hard).}
A pretrained English BERT tokenizer applied to Tunisian Arabic in Latin script (``arabizi'') splits the word \texttt{n7ebek} (``I love you'') as \texttt{[n, \#\#7, \#\#e, \#\#bek]}. Discuss the consequences for a downstream classifier and propose two remediation strategies.

\textit{Solution.}
\begin{itemize}
    \item \textbf{Consequences.} The numeric digit \texttt{7} (which substitutes for the Arabic letter \emph{`}) is treated as foreign; the model sees an unfamiliar pattern with no semantic anchoring. The classifier's representation of the word will be near-random, and its predictions for sentences containing arabizi will be poor.
    \item \textbf{Remediations.}
    \begin{enumerate}
        \item Use a multilingual tokenizer (e.g.\ \texttt{xlm-roberta-base}'s SentencePiece) trained on much broader data including Latin-script Arabic.
        \item Train a domain-specific tokenizer on a Tunisian-Arabic corpus before fine-tuning the model. The merges will then capture frequent dialectal sub-words natively.
    \end{enumerate}
\end{itemize}

\subsection{Python snippet: comparing tokenizers}

\begin{lstlisting}
# pip install nltk transformers
import nltk
from transformers import AutoTokenizer

text = "I've been to Sfax-University; it's amazing!"

# Word-level (NLTK)
nltk.download('punkt_tab', quiet=True)
print("Word tokens (NLTK):", nltk.word_tokenize(text))

# Sub-word (BERT WordPiece)
bert = AutoTokenizer.from_pretrained("bert-base-uncased")
print("BERT sub-word tokens:", bert.tokenize(text))

# Sub-word (GPT-2 byte-level BPE)
gpt2 = AutoTokenizer.from_pretrained("gpt2")
print("GPT-2 BPE tokens:    ", gpt2.tokenize(text))
\end{lstlisting}

The same input produces different tokenisations under different algorithms. Reading the tokenized output is the first debugging step when a model behaves unexpectedly.

\subsection{Exercises}
\begin{enumerate}
    \item For each scenario, recommend a tokenization granularity (word / character / sub-word) and justify briefly:
    \begin{enumerate}
        \item Building a spelling corrector for French SMS;
        \item Training a sentiment classifier on $50$ million Web reviews;
        \item Building a name-entity recognizer for German legal documents.
    \end{enumerate}
    \item Explain in one paragraph why character-level tokenization, despite its $|V| \approx 100$, has \emph{not} replaced sub-word tokenization in modern LLMs.
    \item Tokenize the sentence ``\emph{The CEOs' meeting was rescheduled.}'' with two different policies (whitespace+punctuation, and a sub-word tokenizer of your choice). Comment on the differences for the words \emph{CEOs'} and \emph{rescheduled}.
    \item Carry out three iterations of BPE on the toy corpus $\{\texttt{xay}, \texttt{xay}, \texttt{yax}, \texttt{xy}\}$ where each word has frequency $1$. List the new vocabulary at each step.
    \item (Synthesis.) For a multilingual Tunisian dialect dataset (mixing Arabic script, Latin-script arabizi, and French), discuss the trade-offs between (i) using a single multilingual sub-word tokenizer (e.g.\ XLM-R), (ii) training a domain-specific tokenizer on the dataset, and (iii) using language-specific tokenizers for each script and merging the outputs.
\end{enumerate}

\section{Sparse representations: Bag-of-Words and TF--IDF}
\label{sec:sparse}

\subsection{Motivation}
After tokenization, a text is a sequence of integer IDs. To feed such a sequence to a classical machine-learning algorithm --- logistic regression, $k$-nearest-neighbours, support-vector machine --- we need a \emph{fixed-size vector} representation that does not depend on document length. The simplest such representation is the \textbf{Bag-of-Words} (BoW), which counts how often each vocabulary item appears. BoW and its weighted refinement \textbf{TF--IDF} dominated information retrieval and text classification for two decades; they remain strong baselines, fast to compute, easy to debug, and competitive on small or domain-specific data sets.

\subsection{Approach by problem: are these two reviews similar?}
\textbf{Problem.} Compare the following two hotel reviews:
\begin{itemize}
    \item $d_1$: ``\emph{The room was clean. The room was quiet.}''
    \item $d_2$: ``\emph{The breakfast was great. The room was clean.}''
\end{itemize}
Propose a numerical representation for each that lets us measure their similarity, ignoring word order.

\textbf{Solution sketch.} Build a vocabulary by collecting all distinct words: $V = \{\text{the, room, was, clean, quiet, breakfast, great}\}$ (size $7$). For each document, count how often each vocabulary word appears:
\[ \mathbf{x}_{d_1} = (2, 2, 2, 1, 1, 0, 0)^\top, \quad \mathbf{x}_{d_2} = (2, 1, 2, 1, 0, 1, 1)^\top. \]
The two vectors share many components (e.g.\ both contain ``the'', ``room'', ``was'', ``clean'') and differ on others (\emph{quiet} only in $d_1$, \emph{breakfast / great} only in $d_2$). Their dot product gives a coarse measure of similarity. This embarrassingly simple recipe is the starting point of all classical NLP.

\subsection{Vector space models}

\begin{definition}[Vector space model]
A \textbf{vector space model} (VSM) of a corpus is an embedding $\phi : \mathcal{T} \to \mathbb{R}^d$, where $\mathcal{T}$ is the set of textual objects (words or documents) and $d$ is the embedding dimension. Documents are represented as vectors so that geometric operations (dot product, cosine similarity, distance) translate into linguistic operations (similarity, retrieval, clustering).
\end{definition}

\subsection{Bag-of-Words}

\begin{definition}[Bag-of-Words representation]
Fix a vocabulary $V = \{w_1, w_2, \dots, w_{|V|}\}$. The \textbf{Bag-of-Words} (BoW) representation of a document $d$ is the vector
\[ \phi_{\text{BoW}}(d) = \big(\mathrm{tf}(w_1, d),\ \mathrm{tf}(w_2, d),\ \dots,\ \mathrm{tf}(w_{|V|}, d)\big) \in \mathbb{R}^{|V|}, \]
where $\mathrm{tf}(w, d)$ is the \emph{term frequency} (raw count) of word $w$ in $d$. The representation is named ``bag'' because the order of words is discarded: the document is treated as a multiset.
\end{definition}

\paragraph{Alternative weightings of $\mathrm{tf}$.}
\begin{itemize}
    \item \textbf{Boolean}: $\mathrm{tf}(w, d) \in \{0, 1\}$, indicating presence.
    \item \textbf{Raw count}: $\mathrm{tf}(w, d) = $ number of occurrences.
    \item \textbf{Log-normalized}: $\mathrm{tf}(w, d) = \log(1 + \#(w, d))$, mitigating the bias toward long documents.
    \item \textbf{Length-normalized}: divide raw counts by document length so that all documents have $\ell_1$-norm $1$.
\end{itemize}

\subsection{The document--term matrix}

\begin{definition}[Document--term matrix]
Given a corpus $\{d_1, \dots, d_N\}$ and a vocabulary $V$, the \textbf{document--term matrix} is the matrix $X \in \mathbb{R}^{N \times |V|}$ whose entry $X_{ij} = \mathrm{tf}(w_j, d_i)$.
\end{definition}

In practice $|V|$ is in the tens or hundreds of thousands, while only a few dozen distinct words appear in each document. The matrix is therefore extremely \emph{sparse}: a vast majority of entries are zero. Efficient storage formats (compressed sparse row, CSR; compressed sparse column, CSC) preserve memory by storing only non-zero entries.

\subsection{Limitations of Bag-of-Words}

\begin{itemize}
    \item \textbf{No word order}: the documents ``\emph{Tunisia beat Egypt}'' and ``\emph{Egypt beat Tunisia}'' yield identical BoW vectors despite their opposite meanings.
    \item \textbf{Equal weight to all words}: a function word like ``\emph{the}'' contributes the same magnitude as a domain-specific term like ``\emph{Paracetamol}'', even though the latter is far more discriminative.
    \item \textbf{No semantic similarity}: ``\emph{cat}'' and ``\emph{kitten}'' are orthogonal vectors in $\mathbb{R}^{|V|}$ --- as far apart as ``\emph{cat}'' and ``\emph{democracy}''.
    \item \textbf{Curse of dimensionality}: $|V| \approx 50{,}000$ produces vectors that are computationally cheap (sparse) but statistically noisy --- many features contribute marginally to any specific task.
\end{itemize}

The first two limitations are tackled in this section by TF--IDF; the third requires dense embeddings (Section~\ref{sec:dense}). The first --- the loss of order --- is fundamentally unfixable in a bag-of-words framework and is a key driver of the move to recurrent and Transformer models.

\subsection{TF--IDF}

The intuition is to \emph{down-weight} words that appear in many documents (because they discriminate poorly between documents) and \emph{up-weight} words that appear only in a few (because their presence is informative).

\begin{definition}[Inverse document frequency]
Given a corpus of $N$ documents, the \textbf{document frequency} of a word $w$ is the number of documents that contain $w$ at least once,
\[ \mathrm{df}(w) = \big|\{ d \in \text{corpus} : w \in d \}\big|. \]
The \textbf{inverse document frequency} is
\[ \mathrm{idf}(w) = \log \frac{N}{\mathrm{df}(w)}. \]
Common variants smooth the formula to avoid division by zero and dampen extreme values:
\[ \mathrm{idf}_{\text{smooth}}(w) = \log\left(\frac{N + 1}{\mathrm{df}(w) + 1}\right) + 1. \]
\end{definition}

\begin{definition}[TF--IDF]
The \textbf{TF--IDF} weight of word $w$ in document $d$ is the product
\[ \mathrm{tfidf}(w, d) = \mathrm{tf}(w, d) \times \mathrm{idf}(w). \]
Replacing each entry of the document--term matrix by the corresponding TF--IDF gives the TF--IDF document--term matrix.
\end{definition}

\paragraph{Intuition.}
\begin{itemize}
    \item Words appearing in nearly every document have $\mathrm{df}(w) \approx N$, hence $\mathrm{idf}(w) \approx \log 1 = 0$. They are zeroed out --- the same effect as stop-word removal, but data-driven.
    \item Words appearing in few documents have $\mathrm{df}(w) \ll N$, hence large $\mathrm{idf}(w)$. They retain (and amplify) their term-frequency contribution.
    \item Within a single document, TF--IDF still scales with raw term frequency, preserving emphasis on repeated key terms.
\end{itemize}

\begin{property}[Boundedness of $\mathrm{idf}$]
For a corpus of $N$ documents, $0 \le \mathrm{idf}(w) \le \log N$ when $\mathrm{idf}$ uses the unsmoothed formula. The maximum is reached when $w$ appears in exactly one document, the minimum when $w$ appears in all of them.
\end{property}

\subsection{TF--IDF as a probabilistic measure}

\begin{remark}[Information-theoretic flavour]
$\mathrm{idf}(w) = -\log \frac{\mathrm{df}(w)}{N}$ is the negative log-probability that a randomly chosen document contains the word $w$. A rare word ``surprises'' us when it appears, contributing many bits of information; a common word contributes few. This is a discrete analogue of pointwise mutual information, and the connection has been formalised: TF--IDF approximates a maximum-likelihood estimate under specific generative assumptions.
\end{remark}

\subsection{Worked examples}

\textbf{Example 1 (BoW from scratch --- Easy).}
For the two-document corpus
\begin{itemize}
    \item $d_1$: ``\emph{the cat sat on the mat}''
    \item $d_2$: ``\emph{the cat ate the fish}''
\end{itemize}
build the vocabulary, the document--term matrix with raw counts, and compute $\mathrm{idf}(w)$ for each word.

\textit{Solution.} $V = \{\text{the, cat, sat, on, mat, ate, fish}\}$.
\begin{center}
\begin{tabular}{|c|ccccccc|}
\hline
& the & cat & sat & on & mat & ate & fish \\
\hline
$d_1$ & 2 & 1 & 1 & 1 & 1 & 0 & 0 \\
$d_2$ & 2 & 1 & 0 & 0 & 0 & 1 & 1 \\
\hline
$\mathrm{df}$ & 2 & 2 & 1 & 1 & 1 & 1 & 1 \\
$\mathrm{idf}$ & $\log 1=0$ & $0$ & $\log 2$ & $\log 2$ & $\log 2$ & $\log 2$ & $\log 2$ \\
\hline
\end{tabular}
\end{center}
``\emph{the}'' and ``\emph{cat}'' carry no IDF weight (they appear in every document), while \emph{sat}, \emph{on}, \emph{mat}, \emph{ate}, \emph{fish} all carry weight $\log 2 \approx 0.69$.

\medskip

\textbf{Example 2 (Computing TF--IDF --- Medium).}
A four-document corpus has $\mathrm{df}(\emph{Tunisia}) = 1$ and $\mathrm{df}(\emph{the}) = 4$. Document $d_1$ contains \emph{Tunisia} once and \emph{the} four times.

(a) Compute $\mathrm{tfidf}(\emph{Tunisia}, d_1)$ and $\mathrm{tfidf}(\emph{the}, d_1)$ using the unsmoothed formula.

(b) Comment on the result: which word's contribution dominates the document representation?

\textit{Solution.}
$\mathrm{idf}(\emph{Tunisia}) = \log(4/1) = \log 4 \approx 1.39$. $\mathrm{idf}(\emph{the}) = \log(4/4) = 0$.
$\mathrm{tfidf}(\emph{Tunisia}, d_1) = 1 \times 1.39 = 1.39$. $\mathrm{tfidf}(\emph{the}, d_1) = 4 \times 0 = 0$.
The discriminative word \emph{Tunisia} dominates the representation, even though it appears only once. The function word \emph{the}, which appears four times, is silenced. This is exactly the desired behaviour: a query for \emph{Tunisia} should retrieve $d_1$, while a query for \emph{the} should not.

\medskip

\textbf{Example 3 (Spam filtering with TF--IDF --- Hard).}
Consider a spam-filtering corpus. The word \emph{lottery} appears in $50$ out of $5{,}000$ training emails, all of which are spam. The word \emph{meeting} appears in $1{,}200$ emails, half spam, half legitimate. Argue, using TF--IDF and the discriminative power of each feature, why a logistic-regression classifier built on TF--IDF features will give a much larger weight to \emph{lottery} than to \emph{meeting}.

\textit{Solution.} The IDF of \emph{lottery} is $\log(5000/50) = \log 100 \approx 4.6$, while the IDF of \emph{meeting} is $\log(5000/1200) \approx 1.4$. Already, after the IDF re-weighting, \emph{lottery} contributes more to the representation than \emph{meeting} per occurrence. Logistic regression then assigns weights to maximise log-likelihood; \emph{lottery} is correlated with the spam label with probability $1$, while \emph{meeting} is correlated with the spam label with probability $0.5$. The resulting weight on \emph{lottery} reflects both the higher IDF amplification and the stronger label correlation. The classifier is therefore much more sensitive to the word \emph{lottery} when scoring a new e-mail.

\subsection{Python snippet: from text to TF--IDF in 5 lines}

\begin{lstlisting}
from sklearn.feature_extraction.text import CountVectorizer, TfidfVectorizer

corpus = [
    "the room was clean the room was quiet",
    "the breakfast was great the room was clean",
    "service was slow but the rooms were clean",
]

# Step 1 : raw counts
counts = CountVectorizer().fit_transform(corpus)
print("Raw counts shape:", counts.shape)
print("Raw counts dense:\n", counts.toarray())

# Step 2 : TF-IDF transform
tfidf = TfidfVectorizer().fit(corpus)
X = tfidf.transform(corpus)
print("\nVocabulary:", tfidf.get_feature_names_out())
print("TF-IDF matrix:\n", X.toarray().round(2))

# Inspect the IDF of each token
for term, score in zip(tfidf.get_feature_names_out(), tfidf.idf_):
    print(f"  idf('{term}') = {score:.3f}")
\end{lstlisting}

\subsection{Exercises}
\begin{enumerate}
    \item Build the BoW vocabulary, the document--term matrix, the IDF vector and the full TF--IDF matrix (smoothed) for the corpus
    \begin{itemize}
        \item $d_1$: ``\emph{NLP is fun and useful}''
        \item $d_2$: ``\emph{Tunisia is a country in North Africa}''
        \item $d_3$: ``\emph{NLP is widely used in Tunisia}''
    \end{itemize}
    \item Show that for a corpus of $N$ documents, the maximum possible IDF (unsmoothed) is $\log N$ and is reached when a word appears in exactly one document. Conversely, show that words appearing in every document have $\mathrm{idf}(w) = 0$ and contribute zero TF--IDF weight.
    \item Two documents, $d_1$ and $d_2$, share exactly the same set of words but with different counts. Are their BoW vectors necessarily different? Are their normalised BoW vectors (length-normalised) necessarily different?
    \item TF--IDF is purely lexical and shares no information across morphological variants (\emph{run} and \emph{running} are unrelated columns of the document--term matrix). Propose two engineering remedies (one classical, one modern) and one situation in which each is preferable.
    \item (Open-ended.) Suppose you must build a search engine over a corpus of $100{,}000$ Tunisian-Arabic news articles. Discuss whether you would prefer a TF--IDF index, a dense-vector index (Section~\ref{sec:dense}), or a hybrid of the two. Identify a concrete query that one method handles but not the other.
\end{enumerate}

\end{document}
